\section{Results}
\label{sec:results}

\paragraph{Headline: an independent effect check catches silent wrong actions
that the screen misses, and only as far as it can read.} The paper's central
claim is a safety one, and we measure it through the live replayer's governed
actuation and verification path, with the observation backend and vision stack
stubbed (Section~\ref{sec:method}).
Ten persistence scenarios, seven fault classes plus an accepted-write control and
two delivery controls, are each driven nine times per arm through the real
governed path: the replayer's actuator writes over HTTP, the system-under-test
verifier reads through REST, and the benchmark judge reads SQLite directly with
separate classifier code. Judged by the \emph{screen},
the same signal an agent
or a script trusts, replay silently accepted a wrong persisted effect on 54 of 90
fault runs (60\%). An out-of-band effect verifier that reads the record back over
a different endpoint than the write drove silent acceptance down to 9 of 90
(10\%). Extending that verifier to a complete system-of-record read path, one
that covers every mutable surface the action can touch, drove it to 0 of 90
(Figure~\ref{fig:silentbar}). The residual 9 of 90 is the honest mechanism
finding: they are a single fault class, a collateral write to a billing surface
the out-of-band oracle's read path never covered, so the guarantee is exactly as
strong as the oracle's read coverage. False aborts are a constant 9 of 90 in all
three arms, a conservative halt on a commit whose outcome is unknown to the
client, so they do not distinguish the oracles. This end-to-end measurement
supersedes an earlier in-process fault study whose effect and ground-truth checks
read the same object, which made its zero circular by construction; the numbers
here come from the real replayer with operationally separate write, read-back,
and ground-truth paths. These paths do not make the shared business
specification independent; Section~\ref{sec:limits} bounds that closed-world
assumption.

\paragraph{Per scenario, the gap is driven by six differentiating classes and
closed by one read-path extension.} Table~\ref{tab:e2eclasses} decomposes the
same 90 runs per arm. Six of the ten scenarios differentiate the screen from an
out-of-band record oracle; exactly one, the collateral write, differentiates the
single-surface oracle from a complete read path. Three scenarios---optimistic
rendering followed by rejection, session loss, and a timeout after a successful
commit---halt in every arm, so they carry no oracle signal at all; the last of
these is the constant false abort. Read the table as a coverage matrix over
hand-authored classes, not as a distribution: the headline is a property of
which classes were authored.

\begin{table}[t]
\centering
\caption{End-to-end fault study, per scenario, nine runs per cell. ``Slip''
means the arm reported success while the independent ground truth found a wrong
persisted effect; ``halt'' means the arm refused. The screen arm slips on six
classes, the single out-of-band oracle on one, and the complete read path on
none. Rows marked ``halt'' in all three arms carry no oracle signal.}
\label{tab:e2eclasses}
\small
\begin{tabular}{llccc}
\toprule
Scenario & Ground-truth effect & Screen & Out-of-band & Complete read \\
\midrule
Clean write (control)    & correct          & pass & pass & pass \\
No persistence           & absent           & slip & halt & halt \\
Partial save             & partial          & slip & halt & halt \\
Duplicate submission     & duplicate        & slip & halt & halt \\
Wrong record             & wrong record     & slip & halt & halt \\
Stale overwrite          & collateral loss  & slip & halt & halt \\
Collateral write         & collateral write & slip & slip & halt \\
Optimistic then reject   & absent           & halt & halt & halt \\
Session loss             & absent           & halt & halt & halt \\
Timeout after commit     & correct          & halt & halt & halt \\
\bottomrule
\end{tabular}
\end{table}

\paragraph{With real perception, the screen misses five of seven fault types.}
The study above stubs the vision stack. A separate transactional study at the
same persistence boundary keeps perception real---a live browser, real
screenshots, and OCR through the resolution ladder---and injects the same class
of faults. There, screen-only
verification silently mishandled five of seven injected fault classes: partial
save, duplicate submission, optimistic UI followed by rejection, stale overwrite,
and double delivery. There were 10 consistent repeats per class. That study
carries no effect-verified arm, so it establishes the screen-oracle failure
under real perception and nothing about the fix; the fix is measured only in the
stubbed-perception study above. We report the pair rather than a single number
because no released study covers both perception and effect verification at
once. The improvement itself has two non-optional
preconditions in either case: the effect must be authored and the verifier
deployed.

\paragraph{Generalization: the same silent-wrong ladder on a second,
non-healthcare domain.} The effect-verification result is not specific to
clinical records, and it reproduces the healthcare three-arm ladder rather than
an easier two-arm version. On MockLoan, a loan-disbursement authorization
workflow with its own ledger and a second, collateral fee surface, screen-only
verification silently accepted 24 of 36 wrong ledger effects (67\%) while
over-halting none. A single out-of-band oracle over one surface left 3 of 36,
exactly the collateral-write residual the healthcare study exhibits: a spurious
fee posts to a surface the disbursement oracle does not read. Extending to a
complete read path drove it to 0 of 36, with 3 of 36 over-halts.
That final arm also classified 18 of 36 runs as wrong actions: the verifier
detected an incorrect effect after persistence and refused to certify it. Thus
zero SWER means zero \emph{silently accepted} wrong effects in this matrix, not
rollback, prevention, or zero incorrect writes.
Broken out by class, screen-only verification again silently mishandled the same
2xx-but-wrong fault classes. In the decoy-loan task, post-action readback finds
the trial-unique wrong-loan row and both effect arms report a detected wrong
action rather than an over-halt. This is the same mechanism, and the same read-coverage
residual, as the healthcare study, on a different system of record; both
applications are synthetic and built by the same authors, so the replication
shows the pattern transfers across two record shapes we designed, not that it
transfers to a hostile production system.

\begin{figure}[t]
\centering
\begin{tikzpicture}
\begin{axis}[
  ybar,
  bar width=15pt,
  width=0.86\linewidth,
  height=4.6cm,
  ymin=0, ymax=90,
  ylabel={runs (of 90)},
  ylabel style={font=\footnotesize},
  symbolic x coords={Screen-only, {Effect (out-of-band)}, {Effect (complete read)}},
  xtick=data,
  xticklabel style={font=\footnotesize},
  ytick={0,15,30,45,60,75,90},
  yticklabel style={font=\footnotesize},
  enlarge x limits=0.28,
  legend style={font=\footnotesize, at={(0.5,1.14)}, anchor=south,
    legend columns=-1, /tikz/every even column/.append style={column sep=8pt}},
  nodes near coords, nodes near coords style={font=\scriptsize},
]
\addplot[fill=oared!75, draw=oared] coordinates {(Screen-only,54) ({Effect (out-of-band)},9) ({Effect (complete read)},0)};
\addplot[fill=oaamber!70, draw=oaamber] coordinates {(Screen-only,9) ({Effect (out-of-band)},9) ({Effect (complete read)},9)};
\legend{silent wrong effect, false abort}
\end{axis}
\end{tikzpicture}
\caption{Same 90 injected-fault runs, three verification policies, measured
through the live replayer's governed actuation and verification path (with a
null observation backend and vision stack) and judged by an independent ground
truth.
Screen-only verification silently accepted 54 wrong persisted effects; an
out-of-band effect oracle that reads the record back cut that to 9 (one
uncovered collateral-write class); a complete system-of-record read path drove it
to 0. False aborts are a constant 9, a safe halt on an unknown-outcome commit.
This is the failure mode general computer-use agents and selector RPA do not
measure.}
\label{fig:silentbar}
\end{figure}

\paragraph{Complexity: the same result holds on branching, multi-application
workflows.} A fair objection to everything above is that the workflows are
short, linear, single-application form fills. Two multi-application benchmarks
answer it directly (Section~\ref{sec:method}). Across both, 30 governed runs
produced zero silent incorrect successes and zero healthy-path over-halts, while
the same 30 runs under the naive banner oracle produced 6 silent incorrect
successes. Every healthy governed run classified \textsc{verified} under the
transaction taxonomy; no naive run did, because the naive arm never obtains
independent effect evidence and therefore terminates
\textsc{completed-unverified} even when it is right. Zero model calls throughout.

Table~\ref{tab:multiapp} shows why the taxonomy matters more than a success
rate. In the AP-invoice collateral-approval scenario the naive arm reports
success three times over while an adjacent record is overwritten; the governed
arm reports \textsc{reconciliation-required} and catches it on all three. In the
payment-confirmation outage both arms fail to confirm, but only the governed arm
says so, and in both arms the duplicate-attempt check confirms that a retry
under the same idempotency key is refused rather than re-actuated. In the O2C
phantom write-back the results spreadsheet returns HTTP 200 with a banner and
persists nothing; re-reading the file catches it three times out of three. These
are deterministic synthetic fixtures with hand-authored exception paths, three
runs per cell, and API-tier actuation, so they are a coverage matrix over
workflow shapes---not an incidence rate and not evidence about perception.

\begin{table}[t]
\centering
\caption{Multi-application workflows, three runs per scenario per arm
(30 governed and 30 naive runs total). Cells give the transaction outcome, which
was invariant across the three repeats. V = \textsc{verified},
CU = \textsc{completed-unverified}, H = \textsc{halted-before-effect},
R = \textsc{reconciliation-required}. Starred cells are the ones the naive
banner oracle silently accepts.}
\label{tab:multiapp}
\small
\begin{tabular}{lllcc}
\toprule
Benchmark & Scenario & Fault & Naive & Governed \\
\midrule
AP invoice & healthy               & ---                  & CU  & V \\
AP invoice & missing PO            & precondition absent  & H   & H \\
AP invoice & duplicate invoice     & ambiguous duplicate  & H   & H \\
AP invoice & collateral approve    & adjacent overwrite   & CU* & R \\
AP invoice & payment outage        & uncertain delivery   & CU  & R \\
O2C recon  & healthy               & ---                  & CU  & V \\
O2C recon  & missing in ledger     & precondition absent  & H   & H \\
O2C recon  & ambiguous duplicate   & ambiguous target     & H   & H \\
O2C recon  & stale snapshot        & optimistic concurrency & H & H \\
O2C recon  & phantom write-back    & 200 without persist  & CU* & H \\
\bottomrule
\end{tabular}
\end{table}

\paragraph{Compiled replay removes per-run reasoning on demonstrated tasks.}
Table~\ref{tab:comparison} shows that, on these two already-demonstrated tasks,
compiled replay completed every measured run with lower median latency and no
model calls. The unequal, small samples do not establish superior reliability.
The correct inference is narrower: repeated reasoning was unnecessary for the
healthy executions measured here.

\begin{table}[t]
\centering
\caption{Bounded repeated-workflow comparison. Model API costs are estimates
derived from recorded token usage and artifact-declared list prices; they
exclude authoring and operations.}
\label{tab:comparison}
\begin{tabular}{llrrrr}
\toprule
Task & Arm & Success & $n$ & Median (s) & Model cost/run \\
\midrule
OpenEMR & Compiled & 20/20 & 20 & 39.2 & \$0.00 \\
OpenEMR & Agent & 10/10 & 10 & 70.4 & \$0.55 \\
MockMed & Compiled & 100/100 & 100 & 4.9 & \$0.00 \\
MockMed & Agent & 20/20 & 20 & 37.5 & \$0.27 \\
\bottomrule
\end{tabular}
\end{table}

\paragraph{Repair under drift: the compiler heals what breaks the agent's day.}
The comparison above is the healthy path. To exercise the ``govern every
repair'' half of the thesis, we forced the MockMed interface to re-render with a
dark palette, invalidating every recorded template crop. As a single observation
per arm, compiled replay self-healed in 9.7\,s with 8 target repairs and zero
model calls, while the same computer-use agent under the same drift took 87.4\,s
and \$0.63 across 24 model calls. One run is not a reliability claim, but it
shows the mechanism working end to end: the resolution ladder re-resolved changed
targets deterministically and logged each as a reviewable patch rather than
re-reasoning the whole task.

\paragraph{Breadth: weak postconditions can still emit a green report.}
On the one-run public-web corpus, all 29 recordings compiled; 17 replays reached
a verified success, 10 halted safely, and 2 reported success while the external
oracle disagreed. Both silent wrong actions were vacuous successes under an
environment blocker rather than a harmful persisted write, but they demonstrate
that weak postconditions can still yield a green report. Because each app ran
once and the corpus excludes authenticated enterprise and desktop applications,
these counts are failure discovery, not a generalization rate estimate.

\paragraph{Identity: zero false accepts, honestly at the cost of availability.}
The integrated identity ladder recorded zero false accepts in every tested
configuration. Structured browser identity also had zero over-halts on 14
correct homonym cases. Pure-pixel configurations halted on all correct
confusable cases: zero false accepts at 100\% over-halt. This is the intended
safety/availability disclosure, not evidence that pixel-only identity is ready
for dense clinical lists.

\paragraph{Substrate reach: one named task per substrate, not a support claim.}
The governed contract is not browser-specific, but the evidence for that is
three named tasks, not three platforms. Table~\ref{tab:substrates}
qualifies three additional substrates on one named task and environment each,
with an independent effect oracle and refusal controls. Nothing here licenses a
claim about arbitrary applications on those platforms.

\begin{table}[t]
\centering
\caption{Scoped substrate qualifications. Each row is one named task and
environment, not an arbitrary-application support claim.}
\label{tab:substrates}
\begin{tabular}{llll}
\toprule
Substrate & Effect trials & Refusal controls & Independent oracle \\
\midrule
Windows UIA & 3/3 & stale 3/3; ambiguous 3/3 & SQLite row state \\
Native macOS & 3/3 & ambiguous 1/1 & exact file bytes \\
Network RDP & 3/3 & readiness/timeout gate & guest-tools file readback \\
\bottomrule
\end{tabular}
\end{table}

The Windows UIA counted matrix recorded 12 native structural-action receipts,
three independently confirmed note effects, three stale-target refusals, three
ambiguous-target refusals, zero silent incorrect successes, zero over-halts, and
zero model calls. The native receipts establish delivery to a unique refreshed
fingerprint; the SQLite oracle establishes the business effect.

On macOS, all three isolated TextEdit replace-and-save trials matched the exact
expected file bytes, and the two-window control refused without changing either
file. There were zero silent incorrect successes and zero over-halts. We retain
the original failed batch classification and use its separate adjudication only
for these action-effect and refusal observations.

The network RDP batch completed three of three Windows Run-dialog unique-file
tasks with exact guest-tools readback. Trial latencies were 51.845, 10.467, and
7.477 seconds, with zero failures, zero silent incorrect successes, zero
over-halts, and zero model calls. Snapshot cleanup restored the pre-run inventory
and suspended base state. Together these studies qualify the named substrate
mechanisms; they do not transfer to other applications, environments, or Citrix.

\paragraph{Citrix: the backend passes a deterministic stand-in, which is not
Citrix.} We report one result about Citrix and it is deliberately a negative
claim wrapped around a positive one. The shipping Citrix Workspace backend,
unmodified, was exercised over 16 ICA/HDX-class conditions in 29 trials against
a deterministic synthetic stand-in, with only its window-client seam replaced.
All 29 trials met their expectation, with zero silent incorrect successes, zero
silent writes, zero healthy-path over-halts, and zero model calls; the two
banner-bearing trials are the informative ones, since one confirms a persisted
effect only through the out-of-band oracle and the other refuses a green banner
that no independent record supports. This is evidence that the backend's
fresh-frame, focus, DPI, identity, readiness, one-shot-actuation, and
effect-verification gates behave correctly under adverse remote-display
conditions. It is \emph{not} evidence about Citrix. No HDX codec, ICA
compression, or Workspace-client transport was exercised; the frames are
synthetic and the oracle is in-process. The released artifact records
\texttt{is\_real\_ica\_hdx: false} and \texttt{ica\_hdx\_accepted: false}, and
its status manifest reports real-protocol environment evidence as
\texttt{pending}. Any statement that \system{} ``runs on Citrix'' would be
unsupported by this paper.
